The three agents built across this course represent a clear progression in the trade-off between expressiveness and controllability. The HW1 code-generation agent was the most expressive: given a question, the LLM produced arbitrary Python code that could express essentially any computation. This flexibility came at a cost -- the agent was a black box. If a question produced a wrong answer, it was often unclear whether the fault lay in the LLM's reasoning, a subtle code bug, or an incorrect assumption about the data. I had little recourse but to try to improve the system prompt instructions or the architecture itself, while also trying to keep those general to other questions. The HW2 memory-based agent improved the user experience by maintaining conversational context across turns, but the underlying execution model was unchanged; code generation remained the mechanism, and the same interpretability problems applied.

The HW3 plan-and-execute architecture trades expressiveness for determinism, and that trade-off pays off significantly in reliability and debuggability. Because all computation is expressed as a strictly-typed JSON plan over a fixed operator set, the executor is fully deterministic: the same plan on the same dataset will always produce the same result. This makes failures easy to isolate -- if a question produces a wrong answer, the problem is localized to either the plan (the LLM reasoned incorrectly) or a specific operator (the implementation is wrong). This division between planning and execution meant that debugging was much, much easier. If the planning was poor, I adjusted the instructions or implemented more operators to offer more support. The clearly defined operator interface also enforces a clean separation between reasoning and execution, which means operators can be tested independently and trusted unconditionally once verified. The step schema clearly expressed what the planner wanted to achieve, so I could verify that an operator supported that use properly.

The primary limitation of this design is that the operator set forms a hard ceiling on expressibility. Questions that fall naturally outside the predefined operators -- questions that I have not conceived of -- cannot be answered without first extending the operator library. A code-generation agent handles these cases effortlessly (mostly), since it can synthesize any valid Python. The plan-and-execute agent instead forces the system designer to anticipate the space of needed computations upfront and encode them explicitly. In practice, this constraint is often acceptable: the gain in reliability, interpretability, and reproducibility outweighs the loss in flexibility for well-scoped analytical workloads. But it does mean the agent is less general-purpose than its code-generating predecessors. In completing this assignment, I considered what a hybrid approach might look like: suppose one of the supported operators \textit{was} more flexible -- an LLM code-generate-and-execute call. This could increase flexibility, while still maintaining some level of determinism. Additionally, a developer could monitor what functions were being generated and decide whether to implement those into the primary system.
